\chapter{Discussion}
\label{chap:discussion}

The sensitivity estimates of Section~\ref{sec:sensitivity_comparison} rest on a small set of modelling simplifications that bound their interpretation. The remainder of this chapter addresses the linear-interference omission (Section~\ref{sec:disc_interference}), the absence of systematic uncertainties (Section~\ref{sec:disc_systematics}), and the deferred choice between the three discriminants (Section~\ref{sec:disc_discriminant}).

\section{Interference between the SM and BSM amplitudes}
\label{sec:disc_interference}

The signal samples include only the quadratic EFT term $d^{2}|\mathcal{M}_{8}|^{2}$, while the linear interference term $2d\,\Re(\mathcal{M}_{\mathrm{SM}}^{*}\mathcal{M}_{8})$ is neglected. The quadratic term is positive definite and depends only on the magnitude of the Wilson coefficient, whereas the linear term changes sign with $d$ and must be modelled per operator. The omission is not uniform across the three operator families, since its size tracks the polarisation structure of Chapter~\ref{chap:sm}.

Scalar operators excite only longitudinal modes and therefore share the helicity structure of the SM quartic amplitude, producing the largest linear interference. The quadratic-only model is consequently the least complete on the family with the smallest absolute sensitivity, $Z_A=1.5$ on $S_0$ (Section~\ref{sec:sens_scalar}). Mixed operators lie between the two extremes: their longitudinal component contributes a reduced linear interference, while the transverse component contributes through the quadratic term alone. Tensor operators are orthogonal to the SM amplitude in helicity space, so their linear interference is suppressed. The quadratic-only approximation is therefore most accurate on the family with the largest absolute sensitivity.

\section{Remaining systematic uncertainties}
\label{sec:disc_systematics}

The Asimov significances reported in this thesis include a per-bin $30\%$ Gaussian background normalisation uncertainty through Equation~\ref{eq:zasymptotic}. This single-parameter treatment is a conservative estimate of the experimental uncertainty on the QCD multijet background, which contributes more than 90\% of the events in the training selection. The remaining experimental and theoretical uncertainties are not modelled in the present $Z_A$.

The dominant uncertainties not yet treated are the large-$R$ jet energy and mass scale and resolution, which directly distort the input features the discriminant relies on most heavily, and the modelling of the subdominant $V$+jets, $t\bar{t}$ and diboson backgrounds, whose per-process normalisation and shape uncertainties are not separately propagated. On the signal side, the renormalisation and factorisation scales, the parton distribution functions and the parton-shower matching contribute additional uncertainties that propagate as both normalisation and shape distortions. A realistic measurement would profile each of these effects as a separate nuisance parameter, both relaxing the conservative $30\%$ envelope on the QCD multijet component and introducing the smaller, but non-negligible, contributions from the remaining sources.

The three discriminants are not expected to respond identically to a full profiling. The multi-class architecture, whose softmax outputs distribute probability across the four classes, has its sensitivity concentrated in a narrower region of the discriminant score than the binary models, and may consequently shift relative to the binary architectures once the per-process and shape nuisances are introduced.

\section{Choice of final analysis discriminant}
\label{sec:disc_discriminant}

The three architectures occupy distinct trade-offs in the comparison of Section~\ref{sec:sensitivity_comparison}. The specialised classifier reaches the largest absolute $Z_A$ on the scalar and mixed families at the cost of producing three independent networks, each undefined outside its training pair. The global binary classifier loses $29\%$ to $32\%$ of the specialised significance on the scalar family but matches it within $\pm 5\%$ on the mixed family and overtakes it by $17\%$ to $36\%$ on the tensor family, in exchange for a single network that covers the full nineteen-operator basis. The multi-class classifier loses a further $7\%$ to $21\%$ on the tensor family, $31\%$ to $39\%$ on the mixed family and $25\%$ to $27\%$ on the scalar family relative to the specialised model, in exchange for a per-family discriminant whose interpretation as an operator-family identifier becomes useful after the introduction of linear interference samples in training discussed in Section~\ref{sec:disc_interference}. A definitive choice between the three architectures requires the remaining systematic-uncertainty profiling of Section~\ref{sec:disc_systematics}, which is deferred to the analysis proper. Subject to those studies, the multi-class classifier at the v3 plus equal-$\sigma$ configuration is the preferred architecture of this thesis, since its per-family output becomes a direct operator-family identifier.
